\section{Exact output}\label{sec:exact}

A certified gap below $2^{-q}$ is not yet an exact solution. For rational
quadratics the optimal value and a minimizer have bounded height, so a
small enough certified gap can be turned into an exact answer that is
checked against the original model. This section gives that conversion,
an acceptance rule that needs neither uniqueness nor growth, and the
corresponding statements for polynomial factors.

\subsection{Rational height of quadratic optima}
After substituting fixed coordinates, write the quadratic objective as
\begin{equation}\label{eq:intquad}
F(x)=\frac{x^\T Hx+h^\T x+k}{\Delta},
\end{equation}
where $H$ is a symmetric integer matrix, $h$ and $k$ are integral, and
$\Delta$ is a positive integer such that every box endpoint lies in
$\Delta^{-1}\Z$. (A rational coefficient of $x_ix_j$, $i\ne j$, is split
equally between $H_{ij}$ and $H_{ji}$, so $\Delta$ must absorb a factor
two.) Put
\begin{equation}\label{eq:height}
C_H=\max\bigl\{1,\max_{i,j}\abs{H_{ij}}\bigr\},\qquad
R=\Delta\,(2nC_H)^n,\qquad \Omega=\Delta R^2 .
\end{equation}
These numbers have bit length polynomial in $I$; they are never
enumerated.

\begin{lemma}[Rational height]\label{lem:height}
Problem~\eqref{eq:model} with objective~\eqref{eq:intquad} has a global
minimizer whose coordinates have a common denominator at most $R$, and
$\OPT$ is a rational number with denominator at most $\Omega$. Every global
minimizer that is isolated (in particular a unique one) has coordinates
with a common denominator at most $R$.
\end{lemma}

\begin{proof}
Fix the integer coordinates of some global minimizer; the remaining problem
is a continuous box QP whose minimizers are global minimizers of
\eqref{eq:model}. Among them choose $x$ for which the smallest face of the
continuous box containing $x$ has least dimension. Let $S$ be the set of
continuous coordinates strictly between their bounds at $x$, and $A$ its
complement, so $x_A=t_A/\Delta$ with $t_A$ integral. Since $x$ is a local
minimizer in the relative interior of its face, $\nabla_SF(x)=0$ and
$H_{SS}\succeq0$. If $H_{SS}$ were singular, a nonzero $d\in\R^S$ with
$H_{SS}d=0$ would satisfy $F(x+\tau d)=F(x)$ for all $\tau$ (the first- and
second-order terms vanish), and moving along $d$ until a coordinate of $S$
reaches a bound would give a minimizer on a face of smaller dimension.
Hence $H_{SS}$ is nonsingular, and stationarity reads
\[
(2H_{SS})(\Delta x_S)=-\Delta h_S-2H_{SA}t_A ,
\]
an integral system. By Cramer's rule $x_S$ has a common denominator dividing
$\Delta\det(2H_{SS})$, and
$\abs{\det(2H_{SS})}\le\abs S!\,(2C_H)^{\abs S}\le(2nC_H)^n$. Together with
$x_A\in\Delta^{-1}\Z^A$ this gives a common denominator at most $R$. If $x$
has common denominator $\rho\le R$, then $\Delta\rho^2F(x)$ is an integer,
so $\OPT$ has denominator at most $\Omega$. For an isolated minimizer the
singular case cannot occur, since $x+\tau d$ would be a minimizer for all
small $\tau$, and the same argument applies directly.
\end{proof}

\begin{lemma}[Uniqueness implies growth]\label{lem:unique-growth}
If a quadratic $F$ has a unique minimizer $x^*$ on a bounded mixed-integer
box $X$, then~\eqref{eq:growth} holds for some $g>0$.
\end{lemma}

\begin{proof}
Otherwise there are $x^k\in X\setminus\{x^*\}$ with
$(F(x^k)-\OPT)/\norm{x^k-x^*}^2\to0$. Boundedness gives
$F(x^k)\to\OPT$, so compactness and uniqueness give $x^k\to x^*$, and the
integer coordinates of $x^k$ eventually equal those of $x^*$. Put
$\tau_k=\norm{x^k-x^*}$ and $u^k=(x^k-x^*)/\tau_k$; after passing to a
subsequence, $u^k\to u$ with $\norm u=1$ and $u_i=0$ for $i\in I_Z$. Then
\[
\frac{F(x^k)-\OPT}{\tau_k^2}=\frac{\nabla F(x^*)^\T u^k}{\tau_k}
+\frac12(u^k)^\T\nabla^2F\,u^k .
\]
Since $x^*$ minimizes $F$ over the convex continuous box of its integer
slice, the first term is nonnegative; as the left side tends to zero, the
first term is bounded, so $\nabla F(x^*)^\T u=0$ and
$u^\T\nabla^2Fu\le0$. The direction $u$ lies in the tangent cone of the
polyhedral slice, so $x^*+\tau u$ is feasible for small $\tau>0$, and
$F(x^*+\tau u)-\OPT=\frac{\tau^2}{2}u^\T\nabla^2Fu\le0$. This contradicts
uniqueness.
\end{proof}

The lemma is qualitative; a short input can force an exponentially small
$g$.

\subsection{Exact output under growth}

\begin{theorem}[Exact rational output]\label{thm:exact}
Under the hypotheses of Theorem~\ref{thm:approx} and quadratic
growth~\eqref{eq:growth}, the exact minimizer $x^*$ and the exact value
$\OPT$ can be computed, with a certificate that can be checked
independently, in $f_1(p,\kappa)(I+1)^{C_1}$ bit operations, where $C_1$ is
an absolute constant. Neither $g$ nor $\kappa$ is an input.
\end{theorem}

\begin{proof}
Run CT$(2^{-q})$ for $q=1,2,4,8,\dots$. After each run, with certified
interval $[\beta,U]$ and incumbent $\hat x$, attempt the following
reconstruction. If $U-\beta<1/(4\Omega^2)$, compute by continued fractions the
unique rational of denominator at most $\Omega$ in $[\beta,U]$, if any; two such
rationals differ by at least $1/\Omega^2$. For each coordinate compute the unique
rational of denominator at most $R$ within distance $1/(4R^2)$ of
$\hat x_i$, if any. If all coordinates are reconstructed, accept the
candidate $\tilde x$ only if it lies in $X$ and $F(\tilde x)$ equals the
reconstructed value. An accepted candidate is a global minimizer: its value
is the unique rational of height at most $\Omega$ in an interval that contains
$\OPT$, and $\OPT$ has height at most $\Omega$ by Lemma~\ref{lem:height}.

Once $2^{-q}\le\min\{1/(4\Omega^2),\,g/(32R^4)\}$, growth gives
$\norm{\hat x-x^*}^2\le(F(\hat x)-\OPT)/g\le1/(32R^4)$, so each coordinate
of $x^*$, which has denominator at most $R$ by Lemma~\ref{lem:height}, lies
within $1/(\sqrt{32}R^2)<1/(4R^2)$ of $\hat x_i$; reconstruction returns
$x^*$ and the check succeeds. Since $g\ge L/\kappa$ and $L,R,\Omega$ have
polynomial bit length, the required $q$ is $\poly(I)+O(\log\kappa)$, and
doubling overshoots it by a factor at most two. Theorem~\ref{thm:approx}
and the polynomial cost of continued fractions prove the bound.
\end{proof}

The heights are those of the original problem. Filtering changes the box
endpoints during the search, but it does not change $R$ or $\Omega$.
Arbitrary feasible values do not lie on a fixed lattice when continuous
coordinates are present; only the optimal value does, which is why the
value must be isolated before it is reconstructed.

\subsection{Exact acceptance without uniqueness}
The following rule accepts any candidate, however it was produced.

\begin{proposition}[Exact acceptance]\label{prop:accept}
Let $\beta\le\OPT$ be a certified lower bound and $x\in X$, with
$F(x)=a/W$ in lowest terms. If $F(x)-\beta<1/(\Omega W)$, then $x$ is a global
minimizer.
\end{proposition}

\begin{proof}
If $F(x)>\OPT$, the two rationals have denominators at most $W$ and $\Omega$, so
$F(x)-\OPT\ge1/(\Omega W)$, contradicting $F(x)-\beta<1/(\Omega W)$ and $\beta\le\OPT$.
\end{proof}

The rule uses the denominator of the actual candidate, which may be large;
it does not assume that feasible values lie on a lattice. Combined with a
procedure that proposes rational stationary points on faces suggested by
the incumbent, it gives exact output for every rational box QP, including
problems with a continuum of minimizers.

\begin{theorem}[Finite exact recovery]\label{thm:finite}
For every nonempty rational mixed-integer box QP, the following procedure
terminates with a global minimizer and $\OPT$, certified by
Proposition~\ref{prop:accept}: for $q=1,2,\dots$, run CT$(2^{-q})$, let
$\hat x$ be its incumbent, snap every continuous coordinate within
$\tau=1/(4nR)$ of a box endpoint to that endpoint, fix the integer
coordinates of $\hat x$, solve the linear system ``gradient zero in the
remaining free continuous coordinates'' within the box by exact linear
programming, and apply Proposition~\ref{prop:accept} to its solution. No
growth constant or optimizer is supplied, and no running-time bound is
claimed.
\end{theorem}

\begin{proof}
Let $\mathcal S$ be the set of global minimizers. By
Remark~\ref{rem:termination} the incumbents satisfy $F(\hat x)\to\OPT$,
and compactness gives $\dist(\hat x,\mathcal S)\to0$. Consider an incumbent
$y=\hat x$ within $\tau/2$ of some $s\in\mathcal S$. Since $\tau/2<1/2$,
their integer coordinates agree. Distinct endpoints of a coordinate differ
by at least $1/\Delta>2\tau$, so snapping is unambiguous, and every continuous
coordinate at a bound at $s$ is snapped to that bound. Let $J_0$ be the
continuous coordinates strictly inside their bounds at $s$, and
$J\subseteq J_0$ those not snapped.

Let $P$ be the polytope of points of the continuous box of the integer slice
of $s$ that agree with $s$ outside $J_0$ and satisfy $\nabla_{J_0}F=0$. For
$x,x'\in P$, $d=x-x'$ is supported on $J_0$ and $H_{J_0J_0}d_{J_0}=0$, so
$F(x)=F(x')+\nabla F(x')^\T d+d^\T Hd/\Delta=F(x')$; hence $F\equiv\OPT$ on
$P\ni s$. A vertex of $P$ solves an integral system whose matrix consists of
rows of $2H_{J_0J_0}$ and unit rows for active bounds; expanding along the
unit rows leaves a square, possibly nonprincipal, submatrix of $2H$, whose
determinant is at most $(2nC_H)^n$ in absolute value. So vertices of $P$
have common denominators at most $R$. Let $\phi(x)$ be the sum, over
$i\in J_0\setminus J$, of the distance from $x_i$ to the endpoint to which
$y_i$ was snapped. Then $\phi\ge0$ on $P$, $\phi$ is linear on $P$, and
$\phi(s)\le\abs{J_0}\cdot\frac32\tau<1/R$. Its minimum over $P$ is attained at
a vertex, where it is a nonnegative rational with denominator at most $R$,
hence zero. So some $s'\in P$ satisfies all snaps, and the linear system
solved by the procedure is feasible.

Let $x$ be any solution of that system. Then $x$ agrees with $s'$ outside
$J$, and $\nabla_JF(x)=\nabla_JF(s')=0$, so $d=x-s'$ is supported on $J$
with $H_{JJ}d_J=0$; as before $F(x)=F(s')=\OPT$. An exact simplex method
returns a vertex solution, whose value $\OPT$ has denominator $W\le\Omega$. Once
the certified gap is below $1/\Omega^2\le1/(\Omega W)$, Proposition~\ref{prop:accept}
accepts it.
\end{proof}

The proof shows existence of finite stopping times, not a useful bound on
them; when the optimal set is not a single point, the grids of
Section~\ref{sec:growth} need not shrink around any one minimizer
(Section~\ref{sec:limits}).

\subsection{Explicit polynomial factors}\label{sec:polynomial}
Suppose each factor is an explicit list of rational monomials of degree at
most $d\ge2$. Valid rational curvature bounds $L_i$ are obtained by
differentiating twice in $x_i$, bounding each monomial of $\partial_{ii}F$
over $\bar X$ by interval arithmetic, and summing upper bounds; any sharper
bound may be used if it comes with a certificate whose checking cost is
polynomial in the input, and that certificate is then counted in $I$.

\begin{corollary}[Polynomial factors]\label{cor:poly}
For fixed $d$, under quadratic growth, algorithm CT$(2^{-q})$ returns a
certified $2^{-q}$-approximation in $f_d(p,\kappa)(I+q+1)^{C}$ bit
operations, with $C$ independent of $p$ and $\kappa$. If all coordinates are
integer, it returns an exact minimizer with a certificate in
$f_d(p,\kappa)(I+1)^{C}$ bit operations.
\end{corollary}

\begin{proof}
Sections~\ref{sec:grids} and~\ref{sec:growth} use only
Definition~\ref{def:curvature}, not the quadratic form of $F$. In the bit
analysis, polynomial values at stage $j$ have denominators dividing
$\Delta_FA_j^{d}$, where $\Delta_F$ is a common denominator of the
coefficients, and explicit evaluation costs a polynomial in $I$ and the bit
length of the nodes. For integer coordinates every feasible value lies in
$\Delta_F^{-1}\Z$, so a certified gap below $1/\Delta_F$ at a feasible
incumbent proves it optimal; this needs $q=O(I)$.
\end{proof}

Continuous polynomial problems need a different exact output. The minimizer
of $x^3-6x$ on $[1,2]$ is $\sqrt2$, with value $-4\sqrt2$, so no rational
reconstruction can succeed. Uniqueness also does not imply quadratic growth
for polynomials: $x^4$ on $[-1,1]$ has a unique minimizer and no
$g>0$. Exponents encoded in binary, or factors given as arithmetic circuits,
need a separate evaluation analysis.
